\documentclass[12pt]{article}
\usepackage[utf8]{inputenc}
\usepackage[T1]{fontenc}
\usepackage{hyperref}
\title{Aheu Language Resources for K-12 Curriculum Development}
\author{Compiled by Claude Sonnet 5 (claude-sonnet-5) for Andr\'as Kornai}
\date{}
\begin{document}
\maketitle

\section{Introduction}
Aheu (also called Thavung, Phon Soung or So; ISO 639-3 thm) is a Vietic (Chutic) language of the Austroasiatic family. About 700 to 750 speakers live in three villages of Song Dao district, Sakon Nakhon, Thailand, and more live in Laos. UNESCO rates it as definitely endangered. Documentation is scarce and there are no educational resources.

\section{Primary Educational Resources}
None found.

\section{Linguistic and Reference Materials}
\begin{itemize}
\item Wikipedia entry on Thavung (Aheu): \url{https://en.wikipedia.org/wiki/Thavung_language}
\item Lexical archive entry with a few recordings; it cites the 2008 grammar by Preedaporn Srisakorn, which was not itself seen online: \url{https://cocoon.huma-num.fr/exist/crdo/lexvo.xql?code=thm}
\item Numerals of Aheu (Max Planck EVA): \url{https://lingweb.eva.mpg.de/channumerals/Aheu.htm}
\end{itemize}

\section{Religious and Community Resources}
\begin{itemize}
\item Joshua Project, Aheu in Thailand: \url{https://joshuaproject.net/index.php/people_groups/13129/TH}
\item Joshua Project, Aheu in Laos: \url{https://joshuaproject.net/index.php/people_groups/13129/LA}
\end{itemize}

\section{Digital and Audio Resources}
\begin{itemize}
\item Endangered Languages Project entry: \url{https://endangeredlanguages.com/lang/1214}
\end{itemize}

\section{Curriculum Development Context}
Almost nothing exists for school use. Any programme would begin with community documentation, a practical orthography, probably based on the Thai script, and teaching in Thai alongside it. The 2008 grammar and the recordings are the main starting points.
\end{document}